\documentclass[11pt]{article}
\usepackage[a4paper,margin=1in]{geometry}
\usepackage{fontspec}
\usepackage[boldfont=false]{xeCJK}
\setCJKmainfont{FandolSong-Regular.otf}
\usepackage{amsmath}
\usepackage{amssymb}
\usepackage{array}
\usepackage{booktabs}
\usepackage{graphicx}
\usepackage{adjustbox}
\usepackage{enumitem}
\setlist[itemize]{topsep=2pt,partopsep=0pt,parsep=0pt,itemsep=2pt,leftmargin=1.4em}
\setlist[enumerate]{topsep=2pt,partopsep=0pt,parsep=0pt,itemsep=2pt,leftmargin=1.6em}
\usepackage{parskip}
\usepackage{fvextra}
\fvset{breaklines=true,breakanywhere=true,fontsize=\small}
\setlength{\parindent}{0pt}

\begin{document}

\begin{center}
  {\LARGE\bfseries Government microeconomic intervention (A Level)}\\[0.35em]
  {\large A Level Economics}
\end{center}
\vspace{0.6em}

This topic looks at how governments correct \textbf{market failure} 市场失灵, make incomes fairer, and act in the labour market.

\section*{Correcting market failure}

A government has many tools to push a market towards a better use of resources.

\begin{itemize}
\item an \textbf{indirect tax} 间接税 on a good with a negative externality (like fuel) raises its price and cuts the quantity used.

\item a \textbf{subsidy} 补贴 for a good with a positive externality (like solar panels) lowers its price and raises the quantity.

\item a \textbf{price control} 价格管制 sets a legal maximum or minimum price.

\item \textbf{tradable permits} 可交易许可证 cap total pollution and let firms buy and sell the right to pollute.

\item \textbf{regulation} 管制 uses rules and laws, such as bans, safety standards and pollution limits.

\item \textbf{state provision} 政府提供 means the government supplies a good directly, often free (schools, hospitals, defence).

\item \textbf{nationalisation} 国有化 is taking an industry into government ownership; \textbf{privatisation} 私有化 is selling a state industry to private owners.

\end{itemize}

\subsection*{Competition policy}

A firm with \textbf{monopoly power} 垄断势力 can raise prices and cut output. \textbf{competition policy} 竞争政策 is the set of laws used to protect competition — for example, blocking mergers that create too much power, banning price-fixing, and fining firms that abuse their position.

\subsection*{Government failure}

Intervention can backfire. \textbf{government failure} 政府失灵 is when the government's action leads to a worse use of resources than before — through poor information, high costs, long time lags, or unintended effects (such as a black market caused by a price control).

\section*{Equity and redistribution}

It is important to separate two words that sound alike:

\begin{itemize}
\item \textbf{equity} 公平 means fairness — people getting what is fair, which may not be an equal share.

\item \textbf{equality} 平等 means everyone getting the \textbf{same}.

\end{itemize}

A fair (equitable) outcome is not always an equal one. Most governments worry about too much \textbf{inequality} 不平等.

Remember the difference between the two things being shared:

\begin{itemize}
\item \textbf{income} 收入 is a flow received over time (wages, rent, interest).

\item \textbf{wealth} 财富 is a stock of assets owned (houses, shares, savings).

\end{itemize}

\subsection*{Measuring inequality}

\begin{itemize}
\item the \textbf{Lorenz curve} 洛伦兹曲线 is a graph. It plots the share of total income against the share of the population, from poorest to richest. A straight diagonal line would mean perfect equality. The further the curve bends below the line, the more unequal the country.

\item the \textbf{Gini coefficient} 基尼系数 turns this into one number between 0 and 1. Zero means perfect equality; one means one person has everything. A higher number means more inequality.

\end{itemize}

\subsection*{Policies to redistribute}

\begin{itemize}
\item a \textbf{progressive tax} 累进税 takes a larger share from higher incomes.

\item \textbf{benefits} 福利金 (also called \textbf{transfer payments} 转移支付) give money to the poor, the old and the unemployed.

\item free state services (education, health care) raise the real living standards of the poor.

\end{itemize}

These bring about \textbf{redistribution} 再分配 — moving income from richer to poorer people. Wealth is harder to redistribute than income, because the rich can hide assets or move them abroad, and wealth taxes are hard to collect.

\section*{Labour markets}

A wage is just a price — the price of labour — so we use demand and supply to study the \textbf{labour market} 劳动力市场.

\subsection*{The demand for labour}

The \textbf{demand for labour} 劳动力需求 is a \textbf{derived demand} 派生需求: firms want workers not for their own sake, but to make goods that sell. So labour demand depends on the \textbf{marginal revenue product} 边际收益产品 (MRP) — the extra revenue one more worker brings in. A worker is worth hiring while the MRP is at least as big as the wage.

\subsection*{The supply of labour}

The \textbf{supply of labour} 劳动力供给 to a job rises when the \textbf{wage} 工资 rises, and also depends on training needed, working conditions, and how many people have the right skills.

In a competitive labour market, the wage settles where the demand for labour equals the supply of labour.

\subsection*{When the market is not competitive}

\begin{itemize}
\item a \textbf{monopsony} 买方垄断 is a single, dominant \textbf{buyer} of labour (for example, the only large employer in a town). It can pay a wage below the competitive level and hire fewer workers.

\item a \textbf{trade union} 工会 is an organised group of workers that bargains for higher pay. A union can push the wage above the competitive level, but this may reduce the number of jobs.

\item a \textbf{national minimum wage} 全国最低工资 is a legal lowest wage. Set above the market wage, it raises low pay but may cause some unemployment. Interestingly, in a monopsony it can raise \textbf{both} the wage and employment.

\end{itemize}

\subsection*{Wage differentials and discrimination}

Wages differ between jobs — this is a \textbf{wage differential} 工资差异. The main reasons are differences in MRP (skilled work is more productive), the training and qualifications needed, and the risk or unpleasantness of the work.

Some wage gaps come from \textbf{discrimination} 歧视 — paying people differently because of gender, race or age rather than their productivity. Governments use equal-pay laws and anti-discrimination laws to reduce this.


\end{document}
